% ARCHIVO GENERADO por scripts/generar_capitulos.py. No editar a mano.
\section{Requerimientos de la iteración 3}
Aprobados por el equipo el 2026-11-20 (versión 1.2). Las reglas RN-01 a RN-37 de las iteraciones anteriores siguen vigentes.

\subsection{Respuestas del equipo (2026-11-20)}

\begin{longtable}{>{\raggedright\arraybackslash}p{2.4000000000000004cm}>{\raggedright\arraybackslash}p{10.5cm}>{\raggedright\arraybackslash}p{1.5cm}}
\toprule
\textbf{Pregunta} & \textbf{Respuesta} & \textbf{Dónde quedó}\\
\midrule\endhead
P-39 Duración & La mensualidad vence el \textbf{mismo día del mes siguiente}: si se paga el 18, vence el 18 del otro mes; si se paga el 31 y el mes no tiene 31, vence el \textbf{último día del mes}. & RN-38\\
P-40 Inicio del periodo & La fecha de cobro se mantiene desde que se obtuvo la membresía. Si el dueño paga 4 días tarde, el mes siguiente se paga normal: \textbf{la fecha de cobro no se mueve ni se extiende}. & RN-39\\
P-41 Varios meses & Sí. & RN-40\\
P-42 Descuentos & Puede haber descuentos. \textbf{Solo el administrador} los crea y habilita, por el motivo que sea; una vez habilitado, el \textbf{operador} puede aplicarlo al registrar un pago. & RN-41, RN-42, RF-208\\
P-43 Medio de pago y comprobante & Sí. & RN-43\\
P-44 Pago mal registrado & El pago mal registrado \textbf{queda en un registro} y se registra de nuevo la información corregida. & RN-44, RF-206\\
P-45 Tarifas & Las gestiona el administrador. & RN-45, RF-204\\
P-33 Inactivación & Se aprueba la recomendación: motivo ``Falta de pago'' u ``Otro''; al pagar se reactiva solo lo inactivado por falta de pago. & RN-46, RF-205\\
P-46 Confirmación & Sí: se envía un \textbf{correo simulado} y se genera una \textbf{factura simulada} que se puede \textbf{descargar}. & RN-47\\
P-47 Consultas & Se aprueba la recomendación: mensualidades por vencer en los próximos 5 días e informe de ingresos por mes para el administrador (este último \emph{Could}). & RF-207, RF-209\\
P-48 Cliente que vuelve & Si el vehículo no está dentro, debe pagar para volver a ingresar; si no viene y no paga, no es problema del parqueadero. Vencida no entra (RN-03); si está dentro, aplican los avisos (RN-13 a RN-17). \textbf{No hay membresía nueva}: se mantiene la regla general. & RN-39\\
P-49 Tipo de descuento & El descuento tiene una \textbf{descripción} y un \textbf{porcentaje} que define el administrador. & RN-42, RF-208\\
P-50 Quién corrige & El operador puede corregir un pago durante las \textbf{24 horas} siguientes a su registro; después, solo el administrador. & RN-44, RF-206\\
P-51 Factura & \textbf{Página para imprimir} que el navegador guarda como PDF. & RN-47\\
\bottomrule
\end{longtable}




\subsection{Lo que ya estaba definido}

\begin{longtable}{>{\raggedright\arraybackslash}p{2.4000000000000004cm}>{\raggedright\arraybackslash}p{10.4cm}>{\raggedright\arraybackslash}p{1.6cm}}
\toprule
\textbf{Tema} & \textbf{Definición} & \textbf{Origen}\\
\midrule\endhead
Servicio & Solo por mensualidad, sin cobro por horas. & RN-01\\
Pago por vehículo & Un propietario con varios vehículos paga una mensualidad por cada uno. & RN-10\\
Tarifas & Carro o camioneta 150 000 COP; moto 80 000 COP. & RN-19\\
Entrada & Solo entra un vehículo con la mensualidad vigente. & RN-03\\
Vencida dentro & Puede salir pero no volver a entrar sin pagar; si paga dentro de los 8 días, sale del proceso de retiro. & RN-13, RN-17\\
Quién registra & El operador registra pagos; el administrador puede hacer todo. & RNF-002\\
Trazabilidad & Cada pago guarda el usuario y la fecha y hora; no hay borrado; queda en la bitácora. & RNF-004, RN-07, RN-34\\
Valor histórico & El pago guarda una copia del valor cobrado; si la tarifa cambia, los pagos anteriores no cambian. & DM-04\\
Cambio de propietario & Los pagos siguen asociados al vehículo. & RN-33\\
\bottomrule
\end{longtable}




\subsection{Reglas de negocio propuestas}

\begin{longtable}{>{\raggedright\arraybackslash}p{1.5cm}>{\raggedright\arraybackslash}p{11.4cm}>{\raggedright\arraybackslash}p{1.5cm}}
\toprule
\textbf{ID} & \textbf{Regla} & \textbf{Origen}\\
\midrule\endhead
RN-38 & \textbf{Fecha de cobro.} La membresía de un vehículo empieza el día de su primer pago; ese número de día es su \textbf{día de cobro}. Cada mensualidad vence el día de cobro del mes siguiente; si ese mes no tiene ese día, vence el último día del mes. Ejemplos: 18/11 $\rightarrow$ 18/12; 31/01 $\rightarrow$ 28/02 (o 29/02) $\rightarrow$ 31/03. La mensualidad está vigente hasta el día de vencimiento inclusive. & [U] P-39\\
RN-39 & \textbf{Pagar tarde no mueve la fecha.} Cada pago cubre desde el vencimiento anterior, sin excepción, aunque el vehículo haya estado fuera varios meses. Si vence el 18/12 y paga el 22/12, el pago cubre hasta el 18/01. Mientras no pague, el vehículo no puede entrar (RN-03). El sistema muestra cuántos meses debe pagar para quedar al día. & [U] P-40, P-48\\
RN-40 & Un pago puede cubrir de \textbf{1 a 12 meses} seguidos. & [U] P-41 (el máximo de 12 es [I])\\
RN-41 & \textbf{Valor del pago} = tarifa vigente del tipo de vehículo $\times$ número de meses $-$ descuento aplicado (si hay). El sistema lo calcula; no se escribe a mano. & [U] P-42, RN-19\\
RN-42 & \textbf{Descuentos.} Solo el administrador crea descuentos, con una \textbf{descripción} (el motivo) y un \textbf{porcentaje} entre 1 \% y 100 \%, y los habilita o deshabilita. Un descuento no se edita: si cambia, se deshabilita y se crea otro. Al registrar un pago, el operador puede aplicar \textbf{un} descuento habilitado; el porcentaje se aplica sobre el valor de todos los meses del pago. & [U] P-42, P-49 (rango y ``no se edita'' son [I])\\
RN-43 & Todo pago registra el \textbf{medio de pago} (efectivo, transferencia o tarjeta). & [U] P-43\\
RN-44 & \textbf{Corrección.} Un pago mal registrado no se borra ni se modifica: queda en el registro marcado como corregido, con el motivo obligatorio, y se registra un pago nuevo con los datos correctos. El pago corregido deja de contar para la vigencia y para los ingresos. El operador puede corregir un pago durante las \textbf{24 horas} siguientes a su registro; después, solo el administrador. & [U] P-44, P-50\\
RN-45 & \textbf{Tarifas.} Solo el administrador registra una tarifa nueva por tipo de vehículo, con la fecha desde la que rige (puede ser futura). Las tarifas anteriores no se editan ni se borran. & [U] P-45\\
RN-46 & \textbf{Motivo de inactivación.} Al inactivar un vehículo o un propietario se elige el motivo ``Falta de pago'' u ``Otro'' (con descripción). Al registrar un pago, si el vehículo o su propietario están inactivos \textbf{por falta de pago}, se reactivan automáticamente; si es por ``Otro'', siguen inactivos y el sistema lo advierte. El sistema no inactiva a nadie por su cuenta. & [U] P-33\\
RN-47 & \textbf{Confirmación.} Al registrar un pago, el sistema envía un \textbf{correo simulado} al propietario y deja disponible una \textbf{factura simulada} como página para imprimir (el navegador la guarda como PDF). La factura lleva número consecutivo, fecha, datos ficticios del parqueadero, propietario (nombre y cédula), vehículo, periodo, meses, tarifa, descuento, total, medio de pago y quién lo registró, con la marca ``Simulación – sin validez fiscal''. & [U] P-46, P-51 (los datos de la factura son [I])\\
\bottomrule
\end{longtable}




\subsection{Requerimientos funcionales}

\subsubsection{RF-201 — Registrar pago de mensualidad [E][U]}
\begin{itemize}
  \item \textbf{Actor:} Operador, Administrador.
  \item \textbf{Entradas:} placa, número de meses (1 a 12), medio de pago y, opcional, un descuento habilitado.
  \item \textbf{Proceso:} el sistema muestra el vehículo, su propietario, la tarifa vigente, los meses pendientes, el periodo que cubrirá el pago (RN-38, RN-39) y el valor calculado (RN-41); el usuario confirma; se guarda con el usuario y la fecha y hora; se reactiva lo que corresponda (RN-46); se envía el correo simulado y queda disponible la factura (RN-47); queda en la bitácora.
  \item \textbf{Efecto:} el vehículo queda al día si el pago cubre hasta hoy o después. Si estaba dentro con la mensualidad vencida, sale del proceso de retiro (RN-17).
  \item \textbf{Prioridad:} Must.
\end{itemize}

\subsubsection{RF-202 — Consultar vigencia de la mensualidad [U]}
\begin{itemize}
  \item \textbf{Actor:} Operador, Administrador.
  \item \textbf{Salida:} estado (\textbf{al día}, \textbf{vence hoy} o \textbf{vencida hace N días}), fecha de vencimiento, día de cobro, meses pendientes y último pago. En la ficha del vehículo y al buscar por placa.
  \item \textbf{Prioridad:} Must.
\end{itemize}

\subsubsection{RF-203 — Consultar historial de pagos [I]}
\begin{itemize}
  \item \textbf{Actor:} Operador, Administrador.
  \item \textbf{Salida:} pagos de un vehículo o de todos los vehículos de un propietario: periodo, meses, valor, descuento, medio de pago, fecha, quién lo registró y si fue corregido. Desde cada pago se abre su factura para imprimir.
  \item \textbf{Prioridad:} Must.
\end{itemize}

\subsubsection{RF-204 — Gestionar tarifas [U]}
\begin{itemize}
  \item \textbf{Actor:} Administrador.
  \item \textbf{Proceso:} registrar una tarifa nueva por tipo con su fecha ``vigente desde'' (RN-45). Ver la tarifa actual y el historial.
  \item \textbf{Prioridad:} Must.
\end{itemize}

\subsubsection{RF-205 — Motivo de inactivación y reactivación al pagar [U]}
\begin{itemize}
  \item \textbf{Actor:} Operador, Administrador.
  \item \textbf{Proceso:} RN-46. La ficha muestra el motivo de la inactivación.
  \item \textbf{Prioridad:} Must.
\end{itemize}

\subsubsection{RF-206 — Corregir un pago mal registrado [U]}
\begin{itemize}
  \item \textbf{Actor:} Operador (dentro de las 24 horas siguientes al registro), Administrador (en cualquier momento).
  \item \textbf{Proceso:} se elige el pago, se escribe el motivo y se registran los datos correctos (RN-44). Ambos pagos quedan en el historial y en la bitácora. Si el operador intenta corregir un pago de más de 24 horas, el sistema se lo impide e indica que lo haga el administrador.
  \item \textbf{Prioridad:} Must.
\end{itemize}

\subsubsection{RF-207 — Mensualidades por vencer [U]}
\begin{itemize}
  \item \textbf{Actor:} Operador, Administrador.
  \item \textbf{Salida:} vehículos cuya mensualidad vence en los próximos 5 días, con el propietario y su celular.
  \item \textbf{Prioridad:} Should.
\end{itemize}

\subsubsection{RF-208 — Gestionar descuentos [U]}
\begin{itemize}
  \item \textbf{Actor:} Administrador.
  \item \textbf{Proceso:} crear un descuento con descripción y porcentaje, habilitarlo y deshabilitarlo (RN-42). Ver la lista de descuentos con su estado.
  \item \textbf{Prioridad:} Must.
\end{itemize}

\subsubsection{RF-209 — Informe de ingresos por mes [U]}
\begin{itemize}
  \item \textbf{Actor:} Administrador.
  \item \textbf{Salida:} total recaudado por mes, por tipo de vehículo y por medio de pago, sin contar pagos corregidos.
  \item \textbf{Prioridad:} Could.
\end{itemize}


